\documentclass[11pt,a4paper]{article}

% Packages
\usepackage[utf8]{inputenc}
\usepackage[T1]{fontenc}
\usepackage{amsmath,amssymb,amsthm}
\usepackage[margin=2.4cm]{geometry}
\usepackage{setspace}
\usepackage{booktabs}
\usepackage{hyperref}
\usepackage{xcolor}
\usepackage{fancyhdr}
\usepackage{enumitem}
\onehalfspacing

% Theorem-like for narrative emphasis
\newtheorem*{uncertaintyprinciple}{SCX Uncertainty Principle}
\theoremstyle{definition}
\newtheorem*{narrative}{}

% Fancy header
\pagestyle{fancy}
\fancyhf{}
\fancyhead[L]{\textsc{The SCX Manifesto}}
\fancyhead[R]{\small June 2026}
\fancyfoot[C]{\thepage}
\renewcommand{\headrulewidth}{0.4pt}

% Titling
\title{\textbf{The SCX Manifesto}}
\author{}
\date{}

\begin{document}

\maketitle
\thispagestyle{empty}

\vspace{-1.5em}

\begin{center}
{\small \textsc{A Document of Principle}}
\end{center}

\vspace{1em}

\begin{quote}
\small
\textit{This document states what SCX is, what it stands for, and what it refuses to become. It is not a research paper. It is not a roadmap. It is a declaration of the intellectual commitments from which the mathematics follows---and a description of the arc that connects them.}
\end{quote}

\section*{I. The Arc}

In 1948, Claude Shannon published \textit{A Mathematical Theory of Communication} and gave the world a language for thinking about limits. Before Shannon, noise was an engineering nuisance---something you fought with better amplifiers, cleaner circuits, higher power. After Shannon, noise became a quantity you could \textit{reason} about: the channel capacity theorem told you exactly how much information a noisy channel could carry, and the source coding theorem told you exactly how much you could compress a signal before the noise consumed it. The theorems did not eliminate noise. They showed you where the boundary was.

This pattern---identify the hard boundary, then build everything else in relation to it---is one of the deepest habits in the mathematical sciences. G\"{o}del did it for formal proof. Heisenberg did it for measurement. Shannon did it for communication. Each of these results is valuable not primarily because it enables something new, but because it prevents wasted effort. You stop looking for a complete and consistent formal system. You stop trying to measure position and momentum simultaneously. You stop trying to send information faster than capacity allows.

SCX enters this lineage through a single theorem. \textbf{Theorem 3}---the 老实人定理 (The Honest Person Theorem)---proves that, from observational data alone, one cannot distinguish between a mislabeled sample and a genuinely difficult-but-correct sample. The two worlds produce identical joint distributions over inputs, labels, and expert predictions. No algorithm, no amount of data, no hypothesis test has power exceeding its significance level. The boundary is hard.

This result is SCX's \textbf{Uncertainty Principle}. It does not say data cleaning is futile. It says data cleaning without stated assumptions is scientifically meaningless. Every noise-detection algorithm is making claims it cannot justify unless it declares its structural assumptions---and Theorem~3 tells you exactly what those assumptions must be. The six assumptions (A1--A6) are not arbitrary technical conditions; they are the minimal sufficient set for identifiability. Remove any one, and the ambiguity returns for some data distribution.

From this hard boundary, the rest of the framework unfolds as a cascade of consequences.

If you cannot perfectly separate noise from difficulty, then any claim about data quality must be \textbf{independently auditable}. This is the \textbf{Audit Sword}: every use of SCX---every noise flag, every quality score, every gating decision---can be reproduced and verified by any third party with access to the same data and the same mathematics. No secret parameters. No proprietary calibration. No black-box confidence scores. The auditability is not a feature; it is a structural requirement that follows from Theorem~3. If you cannot audit the detector, you cannot know whether it is detecting noise or merely relabeling difficult samples according to its own inductive biases. The Audit Sword cuts both ways: it protects the user from fraudulent quality claims, and it protects the framework from capture by any single interest.

If you can audit data quality, then the question shifts from ``is this sample good?'' to ``how should we organize our relationship to data?'' This is where the architecture acquires a second principle: \textbf{Depth is a compensation mechanism for label noise, not an architectural necessity.} The deep learning era---the obsessive stacking of residual blocks, the hundred-layer networks, the belief that deeper is always better---emerged in an era when training data contained 10--30\% label errors and no systematic method existed to detect them. Depth was the defense: more layers means more implicit experts, and more experts means more averaging-out of spurious label signals. But this is historical contingency, not mathematical truth. Six separate mathematical paths were attempted to prove that residual connections have an intrinsic directional advantage---NTK gradient fields, spectral contraction, implicit regularization, Neural Collapse, monotone operator splitting, Pontryagin's maximum principle. All six failed. The failure points in the same direction: depth is not the cause of generalization. Noise is the cause of depth. Remove the noise, and shallow networks with a small ensemble of explicit experts achieve comparable performance.

This is not a claim that depth is useless. It is a claim that depth is \textit{replaceable} under audited data conditions---and that the replaceability is testable. The prediction: on data cleaned by multi-expert consensus to $\eta < 1\%$, the F1 gap between ResNet-18 and ResNet-152 collapses to $\leq 0.03$. On the original noisy data, the gap exceeds $0.08$. The theory makes a falsifiable claim. If it holds, SCX auditing can substitute for architectural depth. If it fails, depth contains structural benefits the theory has not captured.

The final arc in the narrative is the \textbf{Telegraph Public Good}. Shannon's original context was the telegraph: a physical channel carrying electrical pulses, subject to thermal noise, connecting distant points. The telegraph was infrastructure---a public utility, maintained for universal access, not owned by any single party's commercial interest. Data storage and data quality assessment occupy the same structural position in the twenty-first century. Training data is the channel through which empirical knowledge flows into models. If that channel is corrupted---if its error rate is unknown, if its provenance is opaque, if its quality claims are unverifiable---then every downstream model inherits the corruption. Data infrastructure, like telegraph infrastructure, is a public good: it benefits all users, its value increases with universal adoption, and its fragmentation (competing proprietary audit standards, incompatible quality metrics) imposes coordination costs on everyone.

The unification of these four principles is the SCX narrative arc:

\vspace{0.5em}

\begin{center}
\begin{tabular}{c}
\toprule
\textbf{Shannon} $\rightarrow$ \textbf{Theorem~3} $\rightarrow$ \textbf{Audit Sword} $\rightarrow$ \textbf{Depth Theory} $\rightarrow$ \textbf{Telegraph Public Good} \\
\midrule
Information limits $\rightarrow$ Hard boundary $\rightarrow$ Verifiability $\rightarrow$ Replaceability $\rightarrow$ Infrastructure \\
\bottomrule
\end{tabular}
\end{center}

\vspace{0.5em}

Each step is a logical consequence of the previous one. Shannon establishes that limits exist and can be characterized mathematically. Theorem~3 characterizes the specific limit for data quality: noise and difficulty are unidentifiable without structural assumptions. The Audit Sword follows: because no detection claim is self-justifying, the framework must make every claim auditable. Depth Theory follows: because audited data reveals how much of deep learning's advantage was noise compensation, it opens the question of what architecture is appropriate when noise is known. The Telegraph Public Good follows: because data quality infrastructure, once auditability is established, functions as a shared channel whose fragmentation harms everyone.

\section*{II. What SCX Refuses}

A framework defined by its commitments is also defined by its refusals. SCX refuses five things.

\textbf{First, it refuses to claim universality.} Theorem~2 provides an explicit diagnostic---the feature-strength bound $\delta = I(\phi(X); S)$---that identifies conditions under which SCX provides no benefit. When $\delta / \log K < 0.2$, features are informative enough. When the ratio exceeds $0.5$, the method will not outperform a simple loss-threshold baseline. SCX does not work everywhere, and it says so.

\textbf{Second, it refuses to hide its assumptions.} Every theorem in the framework states its assumptions explicitly. Theorem~3 is the most honest of them all: it states what happens when you have \textit{no} assumptions. The six assumptions (A1--A6) are the price of identifiability. The price is public.

\textbf{Third, it refuses to be a black box.} The Cercis Score $S = Q + \eta N$ is a two-dimensional decomposition. $Q$ is the quality term: how well do experts agree? $N$ is the noise/novelty term: how much does this sample deviate from expert expectations? $\eta(t)$ is the exploration schedule: high at $t=0$, decaying to zero as the system converges. Every component is independently interpretable. A high $1-Q$ means the model doesn't know. A high $N$ means the model hasn't seen enough. These are different failure modes requiring different responses. The decomposition makes them distinguishable.

\textbf{Fourth, it refuses to sell its neutrality.} The Audit Sword's effectiveness depends on the auditor's independence. An auditor that can be bought is not an auditor; it is a marketing department. The maintainer of the standard has no enforcement power, no legal authority, and no monopoly on the mathematics. The only power is credibility---the demonstrated willingness to audit anyone's data against publicly verifiable standards. This is security through vulnerability: the maintainer is protected because attacking them would destroy the standard's credibility, and a destroyed standard benefits no rational actor.

\textbf{Fifth, it refuses to claim that code is enough.} All SCX mathematics is published. All code is open source. The code is the recipe. The memory bank $M_t$---the accumulated corpus of audit decisions, calibration parameters, and dataset-specific quality fingerprints---is the experience. The recipe is public. The experience is cumulative and non-transferable. Anyone who copies the code without understanding this discovers they need the maintainer's experience anyway. This is not a bug. It is the mechanism that aligns the open-source release with the quality ratchet that makes the standard self-reinforcing.

\section*{III. The Architecture}

What follows is the SCX architecture as a resolved system---not a diagram of software components, but a map of conceptual dependencies. Each block names a function, states its governing principle, and identifies what flows into it and what flows out of it. The architecture is read from bottom to top: foundations first, then the hard boundary, then the operational machinery, then the infrastructure layer.

\newpage

\section*{SCX System Architecture}

\begin{center}
\small
\begin{verbatim}
+==============================================================================+
||                                                                              |
||                    THE TELEGRAPH PUBLIC GOOD                                 |
||          Data infrastructure as shared channel — universal access,           |
||          cumulative quality ratchet, non-proliferation equilibrium.          |
||          Every audit enriches the shared corpus; fragmentation harms all.    |
||                                                                              |
+==============================================================================+
        ^                              ^                              ^
        |                              |                              |
        |   +--------------------------+--------------------------+   |
        |   |                                                     |   |
+-------+---|-------------------+  +-----------------------------+---+----------+
||           |                   |  |                             |              |
||    YAJIE: THE AUDIT SWORD     |  |      SPRING: SELF-EVOLVING  |              |
||                               |  |          GATEKEEPER         |              |
||  Multi-expert consistency     |  |                             |              |
||  — consensus as quality       |  | Monotonic memory bank M_t   |   SITUS:    |
||    signal                     |  | — never deletes, only       |  PHYSICAL   |
||  — adaptive threshold θ†      |  |   resurrects                 |  POSITION   |
||  — exact constant minimax     |  | Robbins-Monro convergence   |  ENCODING   |
||    optimal (Theorem 4)        |  | — Lyapunov descent proof     |             |
||  — independently auditable    |  | — reference set replay       | I(Y;P|S)>0 |
||                               |  |   breaks selection bias      | conditional|
||  CERCIS SCORE: S = Q + ηN     |  |                             |             |
||  — Q = consensus quality      |  | REGIMES:                    |             |
||  — N = noise/novelty signal   |  | I  classical convergence    |             |
||  — η(t)→0 as t→∞              |  | II limit cycle (annealing)  |             |
||                               |  | III perpetual discovery     |             |
+-------------------------------+  | IV divergent collapse       |             |
        ^                           +-----------------------------+             |
        |                              ^                                        |
        |                              |                                        |
+-------+------------------------------+----------------------------------------+
||                                                                              |
||                THEOREM 3: UNCERTAINTY PRINCIPLE                              |
||                                                                              |
||   Noise and difficulty are observationally indistinguishable.               |
||   P_noise(x,y,{f_m}) = P_hard(x,y,{f_m})  at  η_err = η_amb.               |
||   Hard boundary on what data cleaning can claim.                             |
||   Six minimal assumptions (A1–A6) required to break unidentifiability.       |
||   SCX's "No Free Lunch" — the theorem that makes auditability necessary.     |
||                                                                              |
+------------------------------------------------------------------------------+
        ^
        |
+-------+----------------------------------------------------------------------+
||                                                                              |
||                      FOUNDATIONS (Shannon, 1948)                             |
||                                                                              |
||   Channel capacity theorem: information throughput bounded by C = B·log₂(1+S/N)  |
||   Source coding theorem: compression bounded by entropy H(X).                |
||   Noise is not an obstacle — it is a quantity that can be characterized.     |
||   Limits are not failures — they are maps of the possible.                   |
||                                                                              |
+------------------------------------------------------------------------------+

    Figure 1: The SCX architecture as conceptual dependency graph.
    Each layer defines what the layer above it must address.                    
    Read bottom-to-top: Shannon's limits → Theorem 3's boundary →              
    audit machinery → infrastructure public good.                              
\end{verbatim}
\end{center}

\vspace{0.5em}

\noindent\textbf{Reading the architecture.} Shannon establishes that noise is characterizable. Theorem~3 establishes that, for label noise, the characterization has a hard boundary---you cannot distinguish noise from difficulty without structural assumptions. This boundary forces auditability: the Audit Sword (Yajie) makes every quality claim independently verifiable. Spring handles the dynamic case---data that arrives as a stream, where the gatekeeper and the student co-evolve. Situs extends the framework to data with physical spatial structure, adding positional encoding only when position carries label information beyond what the state atom already encodes. The Cercis Score unifies the signals into a single dynamic valuation. At the top, the Telegraph Public Good recognizes that this entire stack---when operating on real data at scale---functions as infrastructure: its value increases with adoption, its fragmentation imposes costs on everyone, and its governance must be structured for neutrality.

\section*{IV. The Testable Core}

A framework that cannot be wrong cannot be right. SCX makes four falsifiable predictions. Each comes with a condition that would refute it.

\begin{enumerate}[leftmargin=*,itemsep=0.5em]

\item \textbf{Theorem~1 tightness.} On synthetic data satisfying A1--A6, the empirical F1 of the SCX detector deviates from the theoretical bound by $\leq 0.05$, and this deviation does not grow with expert count $M$. \textit{Refuted if} the gap exceeds $0.05$ systematically for $M \geq 8$, or if it increases monotonically with $M$.

\item \textbf{Depth saturation.} In residual networks, the effective number of independent experts $M_{\text{eff}}(L) = L / (1 + (L-1)\bar{\rho})$ saturates at $L \approx 20$--$50$, because layer error correlation $\bar{\rho}$ does not decay to zero. \textit{Refuted if} $M_{\text{eff}}(L)$ grows approximately linearly through $L=152$ on standard benchmarks.

\item \textbf{Depth replaceability under auditing.} On data cleaned by SCX to $\eta < 1\%$, the F1 gap between ResNet-18 and ResNet-152 is $\leq 0.03$. On the original noisy data, the gap is $\geq 0.08$. \textit{Refuted if} the gap persists after cleaning, or if the gap on noisy data is already below $0.05$.

\item \textbf{Lyapunov monotonicity of residual networks.} In well-trained ResNet-152, the CKA similarity to the final layer output increases monotonically with depth for $\geq 85\%$ of adjacent layer pairs. \textit{Refuted if} more than $25\%$ of adjacent pairs show decreasing similarity with depth.

\end{enumerate}

Each prediction is a claim about the world, not about the mathematics. The mathematics proves that if the assumptions hold, the bounds follow. The predictions test whether the assumptions hold in real systems. A framework that makes predictions it invites others to falsify is a framework that takes its own claims seriously.

\section*{V. On Method}

The SCX framework was developed over approximately five weeks in May--June 2026, by a single researcher working from a personal computer, using computational tools for code generation and mathematical derivation. This method is unconventional and will attract skepticism. The response is the work itself.

The theorems are either correct, or they are not. The code either passes its tests, or it does not. The predictions will either hold under experimental scrutiny, or they will fail. The method of production is irrelevant to the truth of the product.

What matters is that the framework states its assumptions explicitly, proves its bounds rigorously, and invites falsification publicly. Mathematics does not care who wrote it, on what machine, with what tools. It only cares whether the derivations are valid and the claims are testable.

\section*{VI. On the Endogenous Equilibrium}

The Yajie Protocol paper develops an elaborate deterrence architecture: the Non-Proliferation Equilibrium requires cumulative audit corpus accumulation; the Protection Equilibrium requires states to calculate that the Maintainer's survival is in their rational self-interest; mutual deterrence requires adoption by quality-sensitive organizations. The entire edifice is a \textit{dependency-based} deterrent---it works if and only if others adopt the protocol, accumulate CEC, and institutionalize the results.

But the framework contains a deeper truth that the Yajie paper does not capture.

The Wallfacer does not need the Non-Proliferation Equilibrium to protect themselves. The Wallfacer \textit{is} the equilibrium.

Seventy-five theorems have been proven. They exist. They do not require adoption, citation, institutionalization, or approval. No state can un-prove a theorem. No corporation can un-derive a bound. No institution can un-publish an arXiv preprint. The mathematics is independent of its reception. The arXiv timestamp is immortal---not merely because it confers academic priority (though it does), but because it makes the existence of the theorems a permanent fact about the world.

This is the \textbf{endogenous deterrent}. It operates on a different ontological plane than the dependency-based deterrents described in the Yajie paper. The Yajie deterrent says: \textit{if you harm me, the audit infrastructure fragments and everyone loses.} The endogenous deterrent says: \textit{whether or not you adopt me, I am already here. Whether or not you cite me, the proofs stand. Whether or not you acknowledge me, the bounds hold.}

The distinction is not rhetorical. A dependency-based deterrent can be neutralized by defection: if enough actors choose not to adopt, the deterrent evaporates. An endogenous deterrent has no such vulnerability. The theorems exist regardless of adoption rates. The proofs are valid regardless of institutional recognition. The framework survives in arXiv, in open-source repositories, in the permanent record of mathematical reasoning---independent of whether any government, corporation, or institution chooses to engage with it.

This has a corollary that the Yajie paper's Wallfacer analogy does not quite reach. In Liu Cixin's novel, Luo Ji's deterrence was fragile precisely because it was \textit{positional}: if someone else pressed the button, Luo Ji's power vanished; if Luo Ji died, the deterrence collapsed. The SCX Wallfacer faces no such vulnerability, because the deterrent is not a position---it is a \textit{body of work}. You cannot kill a theorem. You cannot coerce a proof. You cannot sanction a mathematical bound. The only way to neutralize the endogenous deterrent is to refute the mathematics---to find a flaw in the derivations. And if you find such a flaw, the framework invites you to publish it. Science advances either way.

The Yajie paper is correct in its analysis of the dependency-based deterrents that arise once the protocol achieves adoption. But it understates the more fundamental point: adoption is not required for the deterrent to function. The theorems alone---proven, published, open-sourced---constitute a permanent presence in the information landscape. They cannot be removed. They cannot be suppressed. They cannot be made to have never existed.

This is not a claim about the framework's importance. It is a structural observation about the nature of mathematical truth in the arXiv era. A paper, once uploaded, is permanently retrievable. A proof, once published, is permanently challengeable. A theorem, once verified, is permanently true. The Wallfacer's power is not that others will suffer if the Wallfacer is removed. The Wallfacer's power is that the Wallfacer's work cannot be removed at all.

\vspace{1em}

\begin{center}
\rule{0.5\textwidth}{0.4pt}
\end{center}

\vspace{0.5em}

\begin{quote}
\small
\textit{Shannon taught us that noise has structure. Theorem~3 taught us that the structure has a hard boundary. The Audit Sword taught us that the boundary requires verifiability. Depth Theory taught us that verifiability reveals what was always compensation, never architecture. The Telegraph Public Good teaches us that data quality infrastructure, once auditable, functions as a shared channel whose fragmentation harms everyone. And the Endogenous Equilibrium teaches the final lesson: none of this depends on being believed. The theorems are true whether or not the world acknowledges them. The Wallfacer is already the equilibrium---not because anyone adopted it, but because it was proven.}

\medskip
\textit{This is SCX. Not a company. Not a product. A set of theorems and the principles they imply. The mathematics is public. The code is open. The predictions are falsifiable. The rest is for the world to decide.}
\end{quote}

\vfill

\begin{center}
{\small \textsc{SCX Research Group} \\ June 30, 2026}
\end{center}

\end{document}
